\documentclass{article}
\usepackage[margin=1.5in]{geometry}
\usepackage{booktabs}
\usepackage{tabularx}
\usepackage{parskip}
\usepackage[numbers]{natbib}
\usepackage{url}

\setlength{\parindent}{15pt}

\newcounter{quote}
\newcommand{\qcode}{S\refstepcounter{quote}\thequote}

\begin{document}
	
	
	\begin{flushleft}
	{\LARGE \textbf{Replication-Mini-Project: Revisiting the Effects of Work Schedule Uncertainty on Health and Economic Security}} \\
	
	\vspace{20pt}
	
	\uppercase{Lah Hong Wai}, Bauhaus-Universität Weimar \\
	\uppercase{Xavier Julian Theodosius}, Bauhaus-Universität Weimar
	\end{flushleft}
	
	% introduction
	\section{Introduction}
	
	Work schedule uncertainty \textemdash~such as last minute changes, on-call shifts, and last-minute schedule changes \textemdash~has become increasingly common in low- and middle-wage labor markets. Prior research suggests that these problems may have important implications for workers’ health and economic security.
	
	In Improving Health and Economic Security by Reducing Work Schedule Uncertainty \cite{harknett2021}, the authors examine the effects of Seattle's Secure Scheduling ordinance towards workers, by ensuring greater schedule predictability. They found evidence that the law
	had positive effects on workers’ schedule predictability and stability which improved their well-being. They have also investigated causal effects of schedule predictability on well-being outcomes, and discovered that there was a substantial effect.
	
	The purpose of this report is twofold. The original study's main results were first replicated using the authors’ data, but implemented simplified methods, preserving key variables of the original study. In addition to replication, this report extends the original analysis by examining interaction effect between workers' schedule predictability and job seniority, available in the dataset \cite{harknett2021data}. This extension assesses whether worker well-being outcomes differ across seniority.
	
	\section{Study Design}
	
	\subsection{Primary Research Question \& Expectations/Hypotheses}
	
	The original study is motivated by whether precarious working conditions have meaningful consequences for workers' well being. The authors wanted to find if increasing schedule predictability can improve worker well-being. Supported by the ordinance, they desired to explore whether regulations targeting scheduling practices can mitigate this problem.
	
	The paper addresses whether labour policies can influence workers' outcomes, and the relevance of authorities in influencing schedule unpredictability towards economic disparities. The Seattle Secure Scheduling ordinance therefore provided an opportunity to examine these questions by evaluating changes in worker outcomes following a mandated increase in schedule predictability.
	
	The study analysis was led with several testable hypothesis. These hypothesis were not outlined explicitly, but could be extracted from their results. First, they hypothesised that workers covered by the Seattle Secure Scheduling ordinance benefits in the health and economic security as compared to the non-covered workers. Second, they hypothesised that increased work schedule probability is associated with better well-being, such as increased happiness and better sleep quality.
	
	\subsection{Study Type, Design and Conditions}



	% bibliography
	\bibliographystyle{ACM-Reference-Format}
	\bibliography{references}
	
	
\end{document}